\chapter{Cómo leer esta guía}

Esta guía está pensada para alguien recién graduado en Ingeniería Informática: sabes
programar, conoces estructuras de datos, hilos, HTTP y algo de orientación a objetos, pero
\textbf{no sabes Rust} y no has tocado QML. Al terminarla deberías poder abrir cualquier
archivo del proyecto, entender qué hace y por qué está escrito así, y añadir funciones nuevas
sin romper nada.

\section{El recorrido}

\begin{description}[style=nextline, leftmargin=0pt]
  \item[Parte I. Empezar con buen pie] Un tutorial de Rust (capítulo~\ref{cap:rust}) que
  arranca desde cero y se centra en lo que este proyecto usa, con ejercicios que puedes hacer
  en un proyecto de prácticas aparte. Después, QML (capítulo~\ref{cap:qml}), el lenguaje de la
  interfaz. Si ya sabes Rust, salta al capítulo~\ref{cap:qml}.
  \item[Parte II. Cómo está hecho] La arquitectura completa, las herramientas de compilación,
  el puente entre Rust y QML y, sobre todo, un recorrido \emph{archivo por archivo} por todo el
  código: el núcleo en Rust, el código nativo en C++ y Objective-C++, cada componente de la
  interfaz, los formatos de datos y los flujos completos de principio a fin.
  \item[Parte III. Por qué está así] Las decisiones de diseño con sus alternativas
  descartadas, las medidas de rendimiento y memoria que las justifican y las «trampas» que ya
  nos encontramos (y cómo evitarlas).
  \item[Parte IV. Desarrollar por tu cuenta] Recetas paso a paso para los cambios más comunes,
  cómo depurar, y una colección de ejercicios graduados sobre el propio proyecto.
\end{description}

\section{Convenciones}

\begin{itemize}
  \item Los nombres de archivo aparecen así: \archivo{src/player.rs}. Las rutas son relativas a
  la raíz del repositorio.
  \item El código aparece en bloques. Cuando un bloque es un \emph{extracto} del proyecto se
  indica de qué archivo sale; a veces está recortado (con \cod{...}) para quedarnos con lo
  importante. Abre el archivo real al lado: la guía no sustituye al código, lo acompaña.
  \item Los comandos de terminal van en bloques oscuros:
\begin{terminal}
cargo build --release
\end{terminal}
\end{itemize}

Encontrarás cuatro tipos de recuadro:

\begin{idea}
Algo que conviene retener: un concepto que aparece por todo el proyecto.
\end{idea}

\begin{porque}
La razón de una decisión: qué alternativa había y por qué se descartó.
\end{porque}

\begin{trampa}
Un error que ya cometimos (o que es fácil cometer) y cómo evitarlo.
\end{trampa}

\begin{ejercicio}
Algo para hacer tú. Hacer los ejercicios es la forma más rápida de que esto se quede.
\end{ejercicio}

\section{Documentos hermanos}

La guía es el documento largo. En el repositorio hay otros, más cortos, que conviene tener a
mano:

\begin{description}[leftmargin=1.2em]
  \item[\archivo{README.md}] Qué es la aplicación y cómo se usa (para usuarios).
  \item[\archivo{COMPILAR.md}] Compilar e instalar en macOS y CachyOS/Arch.
  \item[\archivo{AGENTS.md}] La chuleta de quien desarrolla: mapa de archivos, comandos,
  argumentos de depuración, convenciones y trampas. Si añades algo, añádelo también ahí.
  \item[\archivo{docs/ARCHITECTURE.md}] Detalle técnico y medidas, en formato corto.
  \item[\archivo{docs/ROADMAP.md}] Qué está hecho y qué falta.
\end{description}

\begin{trampa}[Esta guía también envejece]
El código cambia; si un detalle de la guía no coincide con el archivo, manda el archivo.
Cuando cambies algo importante, actualiza el capítulo correspondiente
(\archivo{docs/guia/capitulos/}) y vuelve a compilar el PDF.
\end{trampa}
